\chapter{Metodologia}
 
O estudo que sustenta este TCC tem como objetivo verificar e estudar os 
impactos do uso de inteligência artificial (IA) no desenvolvimento de software e no papel dos desenvolvedores .  

Com essa finalidade, a pesquisa será exploratória, utilizando fontes 
bibliográficas de duas áreas principais: a de tecnologia e a da saúde. Os artigos da 
área da saúde oferecem embasamento anterior à utilização em massa das 
inteligências artificiais, contribuindo para o referencial teórico do trabalho no que se 
refere a fenômenos como o viés de automação e o avarento cognitivo, estudados 
muito antes da popularização da IA. Esse referencial é desenvolvido em detalhe no 
Capítulo 2.

O método combina a revisão da literatura científica com estudos de caso, 
realizados por meio da produção de código, tanto com ferramentas de inteligência 
artificial quanto manualmente, pelos próprios autores, e da posterior análise de sua 
qualidade. Para a avaliação qualitativa, adota-se uma noção de qualidade baseada 
na norma ISO/IEC 25010, que define características claras de qualidade de código, 
como:
\begin{itemize}
    \item \textbf{Adequação funcional:} O código realiza corretamente o objetivo para o qual foi escrito, considerando "border cases", casos de borda.
    
    \item \textbf{Confiabilidade:} O código comporta-se de maneira correta diante de entradas inesperadas ou erros sobre sistemas não controlados.
    
    \item \textbf{Manutenibilidade:} Legibilidade, modularidade, padrões de código.
    
    \item \textbf{Segurança:} Ausência de vulnerabilidades e de práticas inseguras, especialmente relevantes em código de baixo nível.
\end{itemize}

\section{Composição do corpus}
O corpus da análise utilizará duas linguagens de programação distintas: C e 
Javascript. A escolha se orienta pelos aspectos de segurança, complexidade e popularidade de cada linguagem, C é uma linguagem de baixo nível que exige 
bastante conhecimento de tópicos complexos do programador, como disposição de 
memória, utilização de bibliotecas do sistema e sistemas operacionais, sendo a 
linguagem mais utilizada para programação de sistemas operacionais até hoje. 
Javascript é uma linguagem amplamente utilizada hoje em desenvolvimento de 
sistemas de alto nível como CRUDs e APIs, tem gerenciamento de memória 
automático com garbage collector e é agnóstica de plataforma, funcionando da 
mesma forma em todos os sistemas operacionais, tem tipagem dinâmica e seus 
erros são geralmente relacionados a lógica.  
Dessa forma tem-sê duas linguagens populares e distintas que expõem tipos 
diferentes de defeito: em C, falhas de memória que ainda assim compilam e 
executam; em Javascript, erros sutis de lógica. Essa diferença na facilidade de 
encontrar erros é central para o fenômeno investigado, pois é dessa forma que o 
viés de automação tende a se manifestar, quando o desenvolvedor aceita sem 
questionar código aparentemente correto.

\section{Procedimento de criação do código}

A etapa empírica consiste na geração de código com ferramentas de IA e na 
escrita manual de código pelos próprios autores deste trabalho, que serve como 
base de comparação. Para garantir a reprodutibilidade, serão documentados: os 
modelos utilizados e suas versões, a data de geração, os enunciados (prompts) e os 
parâmetros de geração. 
Serão apresentados os mesmos problemas para serem resolvidos em C e em 
JavaScript, tanto com o código escrito manualmente quanto com o código gerado 
por inteligência artificial, de modo a permitir a comparação entre as duas formas de 
produção.